\documentclass[conference]{IEEEtran}
\IEEEoverridecommandlockouts
% The preceding line is only needed to identify funding in the first footnote. If that is unneeded, please comment it out.
% \usepackage{cite}
\usepackage{amsmath,amssymb,amsfonts}
\usepackage{algorithmic}
\usepackage{graphicx}
\usepackage{textcomp}
\usepackage{xcolor}
 \usepackage{tabularx}
\usepackage{makecell} % for \Xhline
\usepackage{booktabs}
\usepackage{placeins} % \FloatBarrier keeps figures inside their section
\hbadness=10000 \vbadness=10000 % draft: silence harmless underfull-box warnings (remove for the final version)
% \safefig{options}{path}{name shown if missing}: like \includegraphics, but a missing file
% gives a red placeholder box instead of a LaTeX error.
\newcommand{\safefig}[3]{\IfFileExists{#2}{\includegraphics[#1]{#2}}%
  {\fbox{\parbox[c][2.5cm][c]{0.8\linewidth}{\centering\textcolor{red}{[figure not uploaded yet: #3]}}}}}
% \usepackage[slovak]{babel}

\usepackage[style=numeric]{biblatex}
\renewcommand*{\bibfont}{\footnotesize}
\addbibresource{resources.bib}

\def\BibTeX{{\rm B\kern-.05em{\sc i\kern-.025em b}\kern-.08em
    T\kern-.1667em\lower.7ex\hbox{E}\kern-.125emX}}
\begin{document}

\title{Comparison of Machine Learning vs. Neural Network Approaches in Gesture Recognition from an EMG Armband 
% {\footnotesize \textsuperscript{*}Note: Sub-titles are not captured in Xplore and
% should not be used}
% \thanks{Identify applicable funding agency here. If none, delete this.}
}

\author{\IEEEauthorblockN{Dániel Szabó}
\IEEEauthorblockA{\textit{Institute of Artificial Intelligence} \\
\textit{Technical University of Košice}\\
Košice, Slovakia \\
daniel.szabo@student.tuke.sk}
\and
\IEEEauthorblockN{Barbora Novysedláková}
\IEEEauthorblockA{\textit{Institute of Artificial Intelligence} \\
\textit{Technical University of Košice}\\
Košice, Slovakia \\
barbora.novysedlakova@tuke.sk}
% \and
% \IEEEauthorblockN{3\textsuperscript{rd} Given Name Surname}
% \IEEEauthorblockA{\textit{dept. name of organization (of Aff.)} \\
% \textit{name of organization (of Aff.)}\\
% City, Country \\
% email address or ORCID}
% \and
% \IEEEauthorblockN{4\textsuperscript{th} Given Name Surname}
% \IEEEauthorblockA{\textit{dept. name of organization (of Aff.)} \\
% \textit{name of organization (of Aff.)}\\
% City, Country \\
% email address or ORCID}
% \and
% \IEEEauthorblockN{5\textsuperscript{th} Given Name Surname}
% \IEEEauthorblockA{\textit{dept. name of organization (of Aff.)} \\
% \textit{name of organization (of Aff.)}\\
% City, Country \\
% email address or ORCID}
% \and
% \IEEEauthorblockN{6\textsuperscript{th} Given Name Surname}
% \IEEEauthorblockA{\textit{dept. name of organization (of Aff.)} \\
% \textit{name of organization (of Aff.)}\\
% City, Country \\
% email address or ORCID}
}

\maketitle
\begin{abstract}
Low-cost surface electromyography (sEMG) armbands make gesture recognition
accessible, but classifiers that perform well on the people they were trained
on often fail on a new wearer. We compare classical machine learning (ML)
classifiers and neural networks (NN) for gesture recognition from an
eight-channel Myo armband under a single leave-one-person-out protocol. We
recorded a dataset of 211 recordings of six gestures from five people, and
evaluated 16 ML classifiers, including soft-voting ensembles, and 9 neural
architectures, with window-level deduplication and an amplitude-only control
to avoid inflated scores. On an unseen person, the best ML ensemble reaches
77.6\,\% accuracy per window against 72.4\,\% for the best neural network in
the same search run (67.7\,\% averaged over seeds), whereas the network has
about 300 times fewer learned parameters and is 30 times faster per window.
Per-wearer calibration, augmentation and adaptation of the input to the person
were tested as ways to improve generalisation: calibration is worth about four
points, augmentation did not help, and adaptation added a further four points
for the ML models. In informal live use the system seemed to perform better than these
controlled figures suggest, although this was not measured
\textcolor{red}{[TODO: measure live accuracy and replace this sentence with the
result]}.%
% TODO: measure live accuracy (TEST MODEL, new sitting, band re-donned) and
% replace the "not measured" statement here and in the Introduction.
One wearer with a smaller arm was consistently the hardest to recognise;
excluding her recordings from training changed the accuracy for the other
people by less than one point on average, so the difficulty appears to lie in
her own signal, possibly in a loose fit of the band. Both model families fail
on the same gestures and wearer, which suggests that the data, rather than the
classifier, is the limiting factor.
\end{abstract}


\begin{IEEEkeywords}
electromyography, gesture recognition, Myo armband, machine learning, neural networks, cross-subject generalization
\end{IEEEkeywords}



\section{Introduction}
Gesture recognition from surface electromyography (sEMG) is a practical
interface for prosthetics, assistive devices and robot control, because the
muscle signal can be read non-invasively from the forearm. Low-cost armbands
such as the Thalmic Myo, with eight dry electrodes, have made it possible to
collect such data outside the laboratory \cite{TODO_myo_armband}.

A recurring problem is that accuracy does not transfer between people. A
classifier trained and tested on recordings of the same wearer often exceeds
95\,\%, but the same model applied to a person it has never seen performs
considerably worse, because electrode placement, band orientation, skin
condition and muscle fatigue change the signal \cite{TODO_inter_subject}. In
our data the same ML model scores 97.4\,\% on a wearer who is in its training
data (other recordings held out) and 77.6\,\% on one who is not
(Fig.~\ref{fig:gap}).
Reported results are further inflated when overlapping windows from one
recording appear in both the training and the test set, or when each gesture
is recorded in a single take, so that a model can separate classes by
signal amplitude alone instead of by the gesture itself. A subject-independent
evaluation that excludes the tested person from training is therefore the only
meaningful estimate of how a model behaves in use.

A further practical issue is the quality of the electrode contact. The armband
is built for a limited range of arm circumferences, and on thinner arms some
electrodes may not press firmly against the skin. Such a loose connection is
easy to overlook: on the live plot the signal can look plausible, and only on
closer inspection does it become clear that one or more channels carry little
or no usable information. A model trained or evaluated on such recordings
effectively sees fewer than eight channels without this being visible in the
labels or in the overall accuracy, which makes the problem hard to detect
and to separate from genuine differences between gestures.

Two families of models are commonly used for this task. Classical machine
learning (ML) classifiers operate on hand-crafted features and are small,
fast and easy to train on little data. Neural networks (NN) learn directly
from the windowed time series and can exploit the temporal structure of the
signal, but typically need more data and a GPU
\cite{TODO_ml_emg,TODO_dl_emg}. Which family is preferable for an armband
dataset of realistic size is not obvious, and a fair answer requires that both
are trained, tuned and scored under exactly the same protocol.

The contributions of this work are as follows:
\begin{itemize}
  \item a labelled dataset of 211 recordings of six gestures from five
        people, together with an open recording tool that preserves raw data
        and applies per-wearer calibration;
  \item a unified evaluation pipeline with window-level deduplication,
        subject-grouped folds and an amplitude-only control that reveals how
        much of a score comes from channel levels rather than gesture shape;
  \item a comparison of 16 classical ML classifiers, including soft-voting
        ensembles, and 9 neural network architectures, both recurrent and
        non-recurrent, in terms of accuracy on an unseen person, inference time
        and model size;
  \item an ablation of training options aimed at generalisation to a new
        wearer: augmentation (gain, noise, time warping, circular channel
        rotation emulating a rotated band, and channel dropout emulating a
        lifted electrode) and adaptation of the input to the person from a few
        seconds of unlabelled signal. Every option can be switched on or off
        in the settings of the recording tool, and is enabled by default only
        where the ablation showed a gain;
  \item an analysis of a single hard-to-recognise wearer, comparing training
        with and without her recordings and testing band rotation and
        adaptation as explanations.
\end{itemize}
On a person held out from training, the best ML ensemble reaches
77.6\,\% per window against 72.4\,\% for the best neural network, while the
network is about 300 times smaller and 30 times faster per window. Both
families struggle with the same gestures and with the same wearer, which
points to the dataset rather than the model as the limiting factor.

One of the five wearers, who has a smaller arm, is consistently the hardest to
recognise: a model that has never seen her reads only about 60\,\% of her
windows correctly for the ML models and about 40\,\% for the neural networks,
compared with 72--94\,\% (ML) and 66--80\,\% (networks) for the other four people. To test whether her
recordings distort the training set, we repeated the comparison without them.
For the four remaining people, accuracy changed by less than one percentage
point on average (+0.9 for both families), and in different directions for
different people, from +8.0 (Bruno) to $-2.7$ (Betka) points for the ML models. We
therefore found no consistent effect, although with four people we cannot
exclude a small one. The overall score does rise when she is excluded, but only
because her fold, the hardest, no longer enters the average; such figures are
not comparable. Neither a circular shift of the electrodes, emulating a rotated
band, nor leaving her out explained her results, whereas adaptation of the
input to the person raised her accuracy from about 46\,\% to 60\,\% for the ML
models and from 37\,\% to 63\,\% for the GRU. We
suspect that a loose fit of the band on a thinner arm contributes, in line
with the contact problem described above, but this was not tested.

The rest of the paper is organised as follows. Section~\ref{sec:related}
reviews related work and Section~\ref{sec:dataset} describes the dataset.
Sections~\ref{sec:ml} and~\ref{sec:nn} describe the machine learning
algorithms and the neural network, respectively. Section~\ref{sec:experiments}
presents the experiments, Section~\ref{sec:evaluation} evaluates them, and
Section~\ref{sec:conclusion} concludes the paper.



\section{Related Works}
\label{sec:related}
\emph{Classical pipelines.} Gesture recognition from sEMG has traditionally used
windowed, hand-crafted features such as the mean absolute value, waveform length,
zero crossings and slope sign changes, classified with linear discriminant analysis,
support-vector machines or tree ensembles \cite{TODO_ml_emg}. \emph{Neural
networks.} More recent work learns directly from the windowed signal with
convolutional, recurrent and attention-based networks \cite{TODO_dl_emg} and reports
high accuracy within a subject. \emph{Cross-subject generalisation.} The accuracy
falls when a model meets a new wearer, because electrode placement, band orientation
and anatomy change the signal; the usual remedies are calibration or adaptation to
the new wearer and augmentation \cite{TODO_inter_subject}. Reported figures can also
be inflated when overlapping windows of one recording appear in both the training
and the test set, or when each gesture is recorded in a single take, which is why we
remove duplicated windows and add an amplitude-only control.
\emph{Armbands.} The Myo armband \cite{TODO_myo_armband} has eight dry electrodes
and streams either a raw signal or, in the mode we use, a rectified envelope
computed on the device at about 50\,Hz.
\textcolor{red}{[TODO: replace the placeholder citations by the real references and
add the closest related studies with their reported cross-subject accuracies.]}
\section{Dataset}
\label{sec:dataset}
\subsection{Recording and data structure}
The data are multichannel time series. The Myo armband provides eight sEMG
channels, one per dry electrode pod, which we record in its on-device filtered
mode: the signal is rectified and smoothed into an activation envelope, so
every sample is a non-negative integer (nominal range 0--2047) and the stream
has a measured rate of 49.9\,Hz (median; 46.5--51.2\,Hz across recordings);
the radio link delivers the samples in bursts, so they are not evenly spaced
in time.
One recording is therefore a matrix of $T \times 8$ samples with timestamps and
a gesture label; the median recording has 151 samples (3.0\,s).

Six gestures were recorded: rest, fist, open hand, thumb up, wave left and wave
right. Each wearer held every gesture in a relaxed way, at the effort they
would use in practice, during a single continuous take per gesture. The take
was cut afterwards into labelled clips, and all 211 recordings come from such
clips. The median clip is 3.0\,s long and 81\,\% are between 2 and 4\,s
(the shortest is 1.0\,s, the longest 9.5\,s), which gives each wearer between
113 and 154\,s of signal. The dataset contains five wearers with 40--44
recordings each, 10.6 minutes of signal in total (Table~\ref{tab:dataset}).
Each wearer was recorded in one sitting on a single day, so differences
between sittings cannot be separated from differences between wearers. The
classes are approximately balanced; rest is the most frequent with 44
recordings and fist the least frequent with 29. Every recording was checked
automatically for gaps, saturation and a noisy resting baseline; recordings
with warnings (47 with a noisy baseline, 6 with a resting signal that was not
quiet) were flagged but kept. All 211 clips and the 30 longer takes they were
cut from were also read in full for this paper: none contains missing values,
the timestamps never run backwards, the highest count is 1659 (the
nominal maximum is 2047) and no channel is dead.

\begin{table}[!t]
\centering
\caption{Number of recordings per wearer and gesture.}
\label{tab:dataset}
\footnotesize
\setlength{\tabcolsep}{3.5pt}
\begin{tabular}{lrrrrrrr}
\toprule
Wearer & Rest & Fist & Open & Thumb & Wave L & Wave R & Total \\
\midrule
Barbora & 13 & 4 & 6 & 5 & 8 & 7 & 43 \\
Betka   & 7  & 6 & 7 & 7 & 7 & 7 & 41 \\
Bruno   & 8  & 7 & 7 & 6 & 8 & 7 & 43 \\
Dano    & 6  & 7 & 7 & 8 & 8 & 8 & 44 \\
Miso    & 10 & 5 & 8 & 6 & 5 & 6 & 40 \\
\midrule
Total   & 44 & 29 & 35 & 32 & 36 & 35 & 211 \\
\bottomrule
\end{tabular}
\end{table}

The signals of different wearers are not on the same scale. At rest the
channel levels are around 20 for most channels, but for some wearers several
channels sit much higher, for example about 80--110 on four channels for one
wearer, while the activation of a hard contraction reaches several hundred to
over 1500. We therefore put every wearer on a common scale with an affine
calibration derived from that wearer's own unlabelled recordings. The signal is
then resampled onto an even 60\,Hz grid, so that a window always covers the
same duration, and cut into overlapping windows with 75\,\% overlap.
Windows are 1.0\,s long for the ML models and 0.7\,s for the neural networks.
Because clips cut from the same take overlap, windows that cover the same
stretch of signal are kept only once, which leaves 1787 windows for the ML
models and 2873 for the neural networks. The ML windows per gesture are 325
(fist), 318 (open hand), 327 (rest), 223 (thumb up), 319 (wave left) and 275
(wave right); the majority class (rest) accounts for 18\,\% of them. The ML
models do not read the window directly but 50 features computed from it
(Section~\ref{sec:ml}), whereas the networks read the calibrated window.

\subsection{Data visualisation}
\begin{figure*}[!t]
\centering
\includegraphics[width=\textwidth]{figures/Dano_session_01_REST_20261012-164359-976.png}
\caption{Eight sEMG channels (filtered mode, 50\,Hz) over 30\,s of rest for one
wearer, from an earlier recording of the tool that is not part of the final
dataset.}
\label{fig:raw}
\end{figure*}

Figure~\ref{fig:raw} shows 30\,s of the eight channels of one wearer at rest.
Even without a gesture the channels are not flat: each has its own baseline
level, one channel stays clearly higher than the others, and short bursts of
activity appear on most of them. The signal is therefore a mixture of a
channel-dependent baseline and brief bursts, which is why the channels are put
on a common scale before they are used (Section~\ref{sec:dataset}). This
recording comes from the first test session of the tool and its rest is far more
active than that of the clips in the final dataset (Fig.~\ref{fig:profiles}).

Figure~\ref{fig:profiles} shows the final dataset: the recorded counts averaged
over the labelled part of every clip, per wearer, gesture and channel. Rest is
quiet for most wearers (about 20 counts on every channel for Betka and Dano), but
four channels sit at 86--112 for Miso, two at 81--88 for Bruno and one at 76 for
Barbora, which is why a common scale is needed. Every gesture has its own pattern
of strong channels, but the patterns are close: the open hand and the right wave
have a cosine similarity of 0.93--0.98 between their mean profiles for every
wearer. The wearers also differ in how strongly they make a gesture: the thumb up
is on average three times the rest level for Betka and only 1.6 times for Miso,
against about 5 for Barbora and Dano and 11 for Bruno.
\begin{figure*}[!t]
  \centering
  \safefig{width=\textwidth}{figures/channel_profiles.png}{channel\_profiles.png}
  \caption{Mean level of each of the eight channels (recorded counts, colour on a
  logarithmic scale) for every gesture, one panel per wearer.}
  \label{fig:profiles}
\end{figure*}

To see how well the gestures can be separated, Fig.~\ref{fig:pca} projects the
50 features of every ML window to two dimensions by principal component
analysis (PCA), which does not use the labels; the two components explain
61\,\% of the variance. Rest forms a separate cluster, and the two waves
occupy their own regions, while fist, open hand and thumb up overlap. This
agrees with the errors of the classifiers in Section~\ref{sec:evaluation}.
\begin{figure}[!t]
\centering
\includegraphics[width=\columnwidth]{figures/pca_by_gesture.png}
\caption{PCA of the 50 features per window (1787 windows), one gesture
highlighted per panel.}
\label{fig:pca}
\end{figure}

\FloatBarrier

\section{Machine learning algorithms}
\label{sec:ml}
\subsection{Features}
Every ML window is a $60\times 8$ block (1.0\,s at 60\,Hz) of the calibrated
signal. We compare three feature sets. The \emph{level} set has 50 features: for
each channel the mean, median, 90th percentile, standard deviation, the share of
the total activation that falls on the channel and the logarithm of the mean
(48 values), plus the logarithm of the total activation and the ratio of the
strongest channel to the mean channel. The \emph{classic} set has 88 waveform
features, eleven per channel (RMS, mean absolute value, variance, standard
deviation, waveform length, zero crossings, slope sign changes, peak, integrated
EMG, peak-to-peak and skewness), and \emph{both} has 138. The recorded stream is
already rectified, and the classic features subtract each channel's median, which
removes the activation level that identifies a gesture here; the comparison in
Section~\ref{sec:training} shows the consequence.
\subsection{Classifiers}
Sixteen scikit-learn classifiers were compared. Linear and kernel models and
neighbours read standardised features: shrinkage LDA, logistic regression
($C=0.1$ and $1$), $k$-NN ($k=15$ and $31$, distance-weighted), an RBF-SVM ($C=1$, 3,
10 and 30) and two MLPs (one hidden layer of 128 units, and two of 256 and 128, with
$L_2$ penalties $10^{-3}$ and $10^{-2}$, trained with Adam for at most 1200 iterations
and stopped on 10\,\% of their own training fold). The tree models read the features
as they are: histogram gradient boosting (learning rate 0.05, 300 iterations), and
random forests and extra-trees of 500 trees with a minimum leaf size of one or three
windows. All other settings are the scikit-learn defaults; none was tuned on the
held-out wearer.
\subsection{Soft-voting ensembles}
A soft vote averages the class probabilities of its members with equal weights and
answers with the class of the largest mean. The members must be fed identical
features, nothing is refitted, and the vote is scored from the members' stored
out-of-fold probabilities, so any subset can be evaluated without training.

Averaging the probabilities of 2--4 models (every subset of the
eight best models per family, 154 subsets) did not beat the best single model. The
best ML subset reached 79.3\% on the folds it was searched on but 77.0\% when the
choice was repeated without each wearer, against 77.6\% for the best single model; an
unweighted vote that also contained an SVM and LDA scored 69.5\%, because weaker
members outvote the forests.

\section{Neural Network}
\label{sec:nn}
\subsection{NN Architecture}
All nine networks read the 42-step, eight-channel window after the logarithm and a
per-channel standardisation fitted on the training wearers, and end in one linear
layer over the six classes. Most use a width of 32. The \emph{GRU} is a single
recurrent layer with attention pooling over time; the \emph{CNN-GRU} puts a
depthwise-separable convolution in front of it, and the \emph{circular GRU} a
two-dimensional convolution over time and electrode with circular padding along
the electrodes, because the eight pods form a ring. The \emph{LSTM} has two layers
and classifies from its last state. The \emph{CNN} has three convolutional layers
with global average pooling. The \emph{TCN} stacks three depthwise-separable
convolutions with dilations 1, 2 and 4 and attention pooling. The \emph{channel
Transformer} treats the whole window of each electrode as one token (two layers,
four heads), the \emph{InceptionTime-style} network uses three modules with kernel
sizes 3, 7 and 15 and a residual connection, and \emph{TSMixer} alternates
mixing along time and across features. They have between 3.9 and 20.3 thousand
parameters (GRU 4\,455, TCN 3\,943, CNN-GRU 6\,503, InceptionTime-style 6\,854,
TSMixer 9\,066, circular GRU 9\,927, LSTM 14\,022, CNN 18\,246, channel Transformer
20\,326). ROCKET, which was tried once, fits only a linear read-out on random
convolutional kernels (169 thousand parameters). The GRU and TSMixer are described
in Section~\ref{sec:evaluation}.
\subsection{Training}
\label{sec:training}

\textbf{Protocol.} Every model in this paper, classical or neural, is trained and
scored under the same leave-one-person-out protocol: for each of the five wearers
the model is fitted on the windows of the other four and scored on the windows of
the fifth, whom it has never seen. The reported accuracy is the share of correctly
classified windows of the held-out wearer, averaged over the five wearers. Windows
overlap by 75\% and are de-duplicated, so that no stretch of signal is counted twice.

\textbf{Network training.} The networks read a 0.7\,s window after the per-wearer
affine calibration. They are trained with AdamW (learning rate $3\cdot10^{-3}$,
weight decay $10^{-3}$), cross-entropy with label smoothing 0.05, and class weights
inversely proportional to the class frequencies of the training fold. The batch size
is 64, the learning rate follows a cosine decay over a fixed 150 epochs, the hidden
width is 32 and the dropout 0.25. There is no validation split and no early
stopping: the number of epochs is fixed in advance, so the held-out wearer never
influences when training stops. The selected GRU with attention pooling has about
4.5 thousand learned parameters (25\,KB).

\begin{figure*}[!t]
  \centering
  \safefig{width=0.85\textwidth}{figures/zoo_overview.png}{zoo\_overview.png}
  \caption{The model zoo: leave-one-person-out accuracy against model size (log
  scale). Blue: 16 classical classifiers (adaptation to the wearer, three seeds);
  orange: 9 neural networks (two seeds). Ringed markers: the three best models of
  each family, used in the rest of the paper.}
  \label{fig:zoo}
\end{figure*}

\textbf{Model search.} We did not fix the model in advance but searched for it in
three stages, all under the protocol above.
\emph{(i) Classical models and features.} Sixteen classifiers (logistic regression
with two regularisation strengths, shrinkage LDA, $k$-NN with $k=15$ and $31$, an
RBF-SVM with four values of $C$, two MLPs, histogram gradient boosting, and random
forests and extra-trees with 500 trees and two leaf-size settings each) were trained
on 36 feature setups: three signal representations (the signal as recorded, a zero-phase 8\,Hz
low-pass filtered signal, which our code calls band-passed because for a raw 200\,Hz
stream it is a 20--95\,Hz band-pass, and its moving-RMS envelope) $\times$ three feature sets (amplitude
\emph{level} features, classic waveform-shape features, or both) $\times$ four window
lengths (0.4, 0.7, 1.0 and 1.3\,s). The tree ensembles took the whole top 15; level
features were best at every window length except 0.4\,s, where they tie with both
sets (best single model overall 77.4\% against 74.4\% for both feature sets and
66.0\% for the classic ones); and accuracy rose up to
1.0\,s windows and then flattened (Fig.~\ref{fig:windows}, Section~\ref{sec:experiments}). The best setup was the
low-pass filtered signal with level features and 1.0\,s windows.
\emph{(ii) Networks.} Five architectures (GRU, CNN-GRU, TCN, LSTM and CNN) were first
swept at 0.3, 0.4 and 0.5\,s windows; at 0.5\,s the GRU led with 68.3\% ahead of the
CNN-GRU (65.4\%), TCN (64.1\%), LSTM (62.4\%) and CNN (59.9\%). The GRU and CNN-GRU
were then trained at six window lengths from 0.5 to 1.5\,s, and four newer
architectures (circular GRU, InceptionTime-style, channel Transformer and TSMixer)
and ROCKET with a ridge read-out (random convolutional kernels, only the linear
read-out is fitted) were tried on 0.7\,s windows. The GRU was the best network and
peaked at 0.7\,s (72.4\%), but its curve
over window length is jagged and it is a single seed, so we do not read the peak as
a real preference. In the same single-seed run ROCKET reached 69.7\%, with 169
thousand parameters and 8.95\,ms per window against 0.68\,ms for the GRU.
\emph{(iii) The zoo.} All 16 classifiers and 9 networks were then trained again at
their chosen setups (1.0\,s for ML, 0.7\,s for the networks) with several seeds
(three for ML, two for the networks, 150 epochs), and sizes and latencies were
measured (Table~\ref{tab:models}, Fig.~\ref{fig:zoo}). The tree ensembles are the
most accurate (79.4--80.8\%, 5--13\,MB on disk), the best network, the GRU, reaches
67.7\% at 25\,KB, and the other networks lie between 61 and 66\%. The three best
models of each family were carried into the later experiments. Averaged over seeds
the GRU scores 67.7\%, not the 72.4\% of the single search run, and its overall
accuracy moved by about 10 points between seeds in our ablation, so differences of
a few points between networks are within noise.

\textbf{Further training options.} Calibration, adaptation and augmentation, the
soft-voting ensembles and the comparison of training with and without the hardest
wearer are reported in Sections~\ref{sec:ml} and~\ref{sec:experiments}.

\section{Experimental Evaluation}
\label{sec:experiments}
\subsection{Protocol}
The protocol is the one of Section~\ref{sec:training}. We report accuracy per
window as the mean over the five held-out wearers, so that each counts once;
pooling all windows would weight a wearer by the number of windows and gives 79.1\%
instead of 77.6\% for the best ML vote, because the hardest wearer has the fewest.
The spread over wearers is 16\,points (standard deviation) for that vote and is
the main uncertainty of every figure. The per-gesture figures are the share of the
windows of a gesture that are classified correctly. Choosing among hundreds of
models on the same five folds flatters the winner, so for the voting search we also
repeat the whole choice without each wearer (nested check). The amplitude-only
control replaces every window by its eight channel means and repeats training and
scoring; a model that scores the same has used only the activation level.
\subsection{ML models experiments}
Over the 576 distinct fits of 16 classifiers on 36 feature setups, the tree
ensembles hold the whole top 15 and are the best model on all eight setups that use
level features of the recorded or the low-pass filtered signal, by 1.5 to 9.6 points
(77.4\% against 67.8\% for the best other model at the winning setup). On 21 of the
36 setups, mostly weaker ones with classic features or the envelope, a different model
is best, by up to 6 points. The best score of each family is 77.4\% for extra-trees,
76.9\% for random forests, 74.2\% for histogram gradient boosting, 68.4\% for $k$-NN,
67.4\% for an MLP, 66.7\% for the SVM, 64.1\% for logistic regression and 63.9\% for
LDA. Of the signal
representations the low-pass filtered (77.4\%) and the recorded signal (76.9\%) are
equal within noise, and the envelope is clearly worse (64.0\%). The windows give
71.2\% at 0.4\,s, 75.6\% at 0.7\,s, 77.4\% at 1.0\,s and 76.9\% at 1.3\,s. How much
a wearer matters is larger than any of these choices: over the same fits the best
score on the hardest wearer is 56\% and the median 42\%, whereas the easiest wearer
reaches 95\% at best.

\begin{figure*}[!t]
  \centering
  \safefig{width=0.85\textwidth}{figures/window_length.png}{window\_length.png}
  \caption{Accuracy on a held-out wearer against window length. Left: the best ML
  model per feature set. Right: GRU and CNN-GRU (150 epochs, one seed).}
  \label{fig:windows}
\end{figure*}
\subsection{NN experiments}
For the GRU at the six window lengths (one seed each) the accuracy was 66.2\% at
0.5\,s, 72.4\% at 0.7\,s, 65.1\% at 0.9\,s, 68.2\% at 1.1\,s, 64.7\% at 1.3\,s and
65.3\% at 1.5\,s, and for the CNN-GRU 62.9, 60.2, 62.2, 61.6, 58.3 and 63.0\%
(Fig.~\ref{fig:windows}). The weakest of the five folds is only 34--49\% at every length.
With three seeds the GRU baseline scored 72.4, 63.0 and 67.8\% (mean 67.7\%), so the
peak at 0.7\,s is not evidence of a preferred length. Re-running the same seed gave
the identical 72.4\%, so the spread is between seeds and not between runs.
\subsection{ML models and NN comparison}
The two families are compared on the figures that matter for a wearable: accuracy
on an unseen wearer, time per live window and size. In the search run the ML vote
of two forests scored 77.6\% against 72.4\% for the GRU; over three seeds the same
two models score 77.0\% and 67.7\%. The vote needs about 52\,ms per window on one
CPU thread and 17.7\,MB with 1.29 million learned numbers, the GRU 1.8\,ms including
the transfer to the GPU process and 25\,KB with 4\,455 parameters.

Table~\ref{tab:models} compares all models that were trained in terms of accuracy,
size and inference time, and shows what extra help each of them was trained with.

\begin{table}[!htb]
  \caption{All models (leave-one-person-out accuracy per window, \%). Plain = with
  per-wearer calibration only (search stage), Final = with the help listed in the
  header (classical: calibration and adaptation, three seeds; networks: calibration
  only, two seeds). ms = time per live window (classical: one CPU thread, networks:
  GPU). $^\dagger$Single-seed run of the newer-architectures test, not part of the
  zoo; the GRU scored 72.4\% in that run.}
  \label{tab:models}
  \centering
  \footnotesize
  \setlength{\tabcolsep}{2.6pt}
  \begin{tabular}{lrrrrr}
    \hline
    Model & Size & ms & Plain & Final & Gain \\
    \hline
    \multicolumn{6}{l}{\textit{Classical, 1.0\,s windows, calibration + adaptation}} \\
    RF-500-leaf3  & 4.8\,MB & 26.4 & 76.6 & 80.8 & $+4.2$ \\
    RF-500        & 5.7\,MB & 26.5 & 76.6 & 80.5 & $+3.9$ \\
    ET-500-leaf3  & 11.5\,MB & 26.9 & 75.6 & 79.4 & $+3.8$ \\
    ET-500        & 13.0\,MB & 26.6 & 77.4 & 79.4 & $+2.0$ \\
    SVM-C1        & 219\,KB & 2.47 & 65.2 & 77.2 & $+12.0$ \\
    SVM-C3        & 167\,KB & 2.27 & 64.9 & 76.4 & $+11.5$ \\
    MLP-256-128   & 564\,KB & 3.53 & 65.0 & 75.7 & $+10.7$ \\
    MLP-128       & 102\,KB & 2.34 & 67.3 & 74.9 & $+7.6$ \\
    HGB           & 2.9\,MB & 28.6 & 67.8 & 74.6 & $+6.8$ \\
    LogReg-C0.1   & 10.7\,KB & 2.25 & 61.5 & 74.4 & $+12.9$ \\
    SVM-C10       & 157\,KB & 2.33 & 64.2 & 74.4 & $+10.2$ \\
    SVM-C30       & 150\,KB & 4.67 & 64.1 & 73.8 & $+9.7$ \\
    LDA           & 31\,KB & 2.60 & 63.1 & 73.4 & $+10.3$ \\
    $k$-NN-31     & 371\,KB & 5.76 & 66.5 & 72.9 & $+6.4$ \\
    $k$-NN-15     & 371\,KB & 3.55 & 66.3 & 72.8 & $+6.5$ \\
    LogReg-C1     & 10.7\,KB & 2.17 & 62.6 & 69.7 & $+7.1$ \\
    \hline
    \multicolumn{6}{l}{\textit{Networks, 0.7\,s windows, calibration only}} \\
    ROCKET$^\dagger$ & -- & 8.95 & -- & 69.7 & -- \\
    GRU           & 23.5\,KB & 0.81 & -- & 67.7 & -- \\
    TSMixer       & 45.9\,KB & 1.81 & -- & 66.3 & -- \\
    Chan.\ Transf. & 92.2\,KB & 1.25 & -- & 66.0 & -- \\
    LSTM          & 61.3\,KB & 0.52 & -- & 63.9 & -- \\
    Circ.\ GRU    & 47.1\,KB & 1.28 & -- & 63.4 & -- \\
    TCN           & 30.2\,KB & 2.08 & -- & 62.8 & -- \\
    CNN           & 82.6\,KB & 1.09 & -- & 62.6 & -- \\
    CNN-GRU       & 34.9\,KB & 1.49 & -- & 62.3 & -- \\
    Inception-S   & 42.3\,KB & 2.62 & -- & 61.1 & -- \\
    \hline
  \end{tabular}
\end{table}
\subsection{Calibration, augmentation and adaptation}
Three kinds of help for a new wearer were compared, each against a baseline on
identical folds and seeds.

\begin{figure*}[!t]
  \centering
  \safefig{width=0.85\textwidth}{figures/exclusion_per_person.png}{exclusion\_per\_person.png}
  \caption{Accuracy on each remaining wearer for models trained with (light) and
  without (dark) Barbora's recordings; mean of the three recommended models per
  family, dots are single models.}
  \label{fig:barbora}
\end{figure*}
Three kinds of help are available
(Tables~\ref{tab:models} and~\ref{tab:help}). \emph{Calibration} puts every wearer on a
common scale from their own unlabelled signal; switching it off costs the ML models
4.3 points. \emph{Adaptation} standardises the input with statistics of a random 10\%
of the wearer's unlabelled windows (about 9\,s) and adds a further 4.3 points to the ML
pair, so that it rises from 72.7\% without any help to 77.0\% with calibration and
81.3\% with both. It helps every classical model, by 2--4 points for the forests and
10--13 for the SVMs, LDA and logistic regression with $C=0.1$, so the best non-forest
model (SVM) is 3.6 points behind the best forest instead of about 10. The gain
is uneven across wearers (for the ML pair Barbora $+14.6$, Miso $+8.4$, Dano $+5.3$,
Betka $-0.6$ and Bruno $-6.3$ points), and by our rule it passes only on the
robustness probes and then narrowly, so it is the weakest of the defaults we use. For the GRU it
gave $+4.5$ points but with a baseline seed spread of 4.7 (against its own spread
of 1.5 it would have passed, so this is a borderline call), so we left it off; the ML models are
therefore trained with it and the networks without, which slightly favours the ML
models. The gain is uneven: for the GRU it lifted Barbora from 37\% to 63\% but
lowered Betka by 7 points (86\% to 79\%). A classical model trained with adaptation
and used without the wearer's statistics reads 78.2\% against 78.5\% with them
(seed 0), so it can be used before a calibration signal exists. Five \emph{augmentations} (gain, noise, time warp, band rotation, i.e.\ a
circular shift of the eight channels, and a dead electrode) were also tried. Judged by
a fixed rule (a gain of at least one point and twice the larger of the two seed spreads, in at least four
of the five wearers), none is used: gain, rotation and dead electrode hurt, noise and
time warp do nothing. Calibration is uneven across wearers too. Switching it off lowers Bruno from 80\% to
57\% but raises Betka from 83\% to 90\% and Barbora from 46\% to 53\%, and a single
global gain, which keeps the profile across channels intact, raises Barbora from 46\%
to 75\% and lowers Miso from 88\% to 71\%. The mean favours the affine calibration we
use (global gain $+3.6$ points but in only three of five wearers and within noise),
yet the global gain is the most interesting lead for the hardest wearer and
deserves a test with more wearers. Rotation raises the accuracy with a band slid by one electrode by
about 17 points (ML 43.8\%\,$\to$\,61.1\%, GRU 45.6\%\,$\to$\,62.2\%) but costs 4--7
points on unshifted data.

\begin{table}[!htb]
  \caption{What each extra help achieved (accuracy change in points, five wearers,
  three seeds). Reference: ML = soft vote of ET-500 and RF-500-leaf3, GRU = network,
  both with calibration only. Clean = normal test data, shifted = band slid by one
  electrode. The seed spread of the reference is 0.5 points (ML) and 4.7 (GRU), so
  GRU changes below about 9 points cannot be told from noise. n/m = not measured.
  Used = whether ML / GRU training uses it.}
  \label{tab:help}
  \centering
  \footnotesize
  \setlength{\tabcolsep}{2.6pt}
  \begin{tabular}{lrrrrl}
    \hline
    & \multicolumn{2}{c}{ML vote} & \multicolumn{2}{c}{GRU} & \\
    Extra help & clean & shifted & clean & shifted & Used \\
    \hline
    Reference, accuracy \% & 77.0 & 43.8 & 67.7 & 45.6 & -- \\
    Calibration off       & $-4.3$ & $+6.0$ & n/m & n/m & on / on \\
    Global gain only      & $+3.6$ & $+1.1$ & n/m & n/m & off \\
    Per-channel only      & $-8.1$ & $-4.2$ & n/m & n/m & off \\
    Adaptation            & $+4.3$ & $+0.6$ & $+4.5$ & $+3.2$ & on / off \\
    Aug.: gain            & $-4.7$ & $-2.4$ & $-1.8$ & $-0.3$ & off \\
    Aug.: noise           & $+0.0$ & $-1.1$ & $+0.4$ & $+0.8$ & off \\
    Aug.: time warp       & $-0.1$ & $-0.7$ & $-0.4$ & $+0.3$ & off \\
    Aug.: band rotation   & $-6.6$ & $+17.3$ & $-4.2$ & $+16.6$ & off \\
    Aug.: dead electrode  & $-4.2$ & $-1.5$ & $-1.1$ & $+4.3$ & off \\
    Barbora left out      & $+0.9$ & -- & $+0.9$ & -- & off \\
    \hline
  \end{tabular}
\end{table}
\subsection{The hardest wearer}
One wearer is read far worse than the others. We looked for the cause and for a
remedy.


Barbora is the hardest wearer: models
that have never seen her read her at only about 60\% (ML) and 39\% (networks), far
below everybody else. Her errors are not blurred but substituted: with the plain ML
vote 81\% of her fist windows are answered as open hand, all of her open-hand windows
as right wave and 83\% of her thumb-up windows as open hand, while rest and both waves
are read correctly (94--98\%). The model is also as sure of these substituted answers
as of correct ones (mean top probability 0.60 when right and 0.61 when wrong), whereas
for the other four wearers wrong answers are clearly less certain (0.38--0.53 against
0.64--0.82), so the confidence cannot be used to catch her errors. Rotation augmentation did not help (ML 45.9\% $\to$ 41.6\%;
for the GRU 36.7\% $\to$ 56.1\%, but with very different results across seeds and
a loss of 25 points on Miso). We therefore also trained the six recommended models
again with identical settings but without her recordings, and scored them on the
four remaining wearers (Betka, Bruno, Dano, Miso). For a kept wearer both numbers
come from a model that never saw him or her, so the difference is exactly what
Barbora's recordings did for that person.

The effect is small and not consistent (Fig.~\ref{fig:barbora},
Table~\ref{tab:barbora-models}). Averaged over the four wearers, leaving her out
changes the ML models from 85.4\% to 86.3\% and the networks from 73.5\% to 74.4\%
(+0.9 points in both cases; +0.2 to +2.2 for single models). Single-wearer changes
range from $-5.7$ to $+9.5$ points and often change sign between models for the
same wearer (Betka, Miso). With four wearers and two or three seeds this cannot be
separated from noise. The headline scores without her are higher
(86.2--86.4\% for ML and 73.6--75.4\% for the networks) only because the hardest
wearer is no longer in the average, so they are not comparable with the five-wearer
scores. Her own result is the clear loser: without her recordings she is read at
59.7\% (ML) and 39.7\% (networks); the case where the model has seen her recordings
and she wears the band in a new sitting was not measured. Leaving a wearer out is
therefore offered as a setting, off by default.

\begin{table}[!htb]
  \caption{Change in accuracy (points) of the six recommended models on each
  remaining wearer when Barbora's recordings are left out of training.}
  \label{tab:barbora-models}
  \centering
  \footnotesize
  \setlength{\tabcolsep}{3.5pt}
  \begin{tabular}{lrrrrr}
    \hline
    Model & Betka & Bruno & Dano & Miso & Mean \\
    \hline
    ET-500-leaf3       & $-0.4$ & $+5.3$ & $+1.0$ & $+2.8$ & $+2.2$ \\
    RF-500-leaf3       & $-3.2$ & $+9.1$ & $-0.3$ & $-4.2$ & $+0.4$ \\
    RF-500             & $-4.6$ & $+9.5$ & $-0.7$ & $-3.6$ & $+0.2$ \\
    \hline
    Chan.\ Transformer & $-4.1$ & $+4.7$ & $+1.3$ & $+2.6$ & $+1.1$ \\
    TSMixer            & $-2.6$ & $+0.3$ & $-1.5$ & $+8.1$ & $+1.1$ \\
    GRU                & $+6.9$ & $+1.2$ & $-0.1$ & $-5.7$ & $+0.6$ \\
    \hline
  \end{tabular}
\end{table}

\section{Evaluation of Experiments}
\label{sec:evaluation}
\subsection{Errors per gesture and per person}
\begin{figure*}[!t]
  \centering
  \safefig{width=0.8\textwidth}{figures/confusion_both.png}{confusion\_both.png}
  \caption{Confusion matrices with the wearer held out, ML vote (left) and GRU
  (right). Rows are the gesture that was made, columns the answer of the model; the
  large number is the share of the row.}
  \label{fig:confusion}
\end{figure*}
The errors are concentrated in a few gestures and a few wearers. The waves and rest
are easy for both families, whereas fist, open hand and thumb up are where both lose
their points (Table~\ref{tab:gestures}). The three use overlapping forearm muscles, and
on an eight-electrode band their patterns lie close together: in a projection of all
training windows that never sees the labels, the waves and rest form their own
islands while fist, open hand and thumb up overlap in the middle (Fig.~\ref{fig:pca}).
The open hand and the right wave are almost the same signal (cosine similarity of the
mean channel profiles 0.98 in the recorded counts and 0.92 after calibration, against
0.56 between the two waves). Choosing another model therefore removes few of these
errors: of the 1787 windows, 896 (50\%) are classified correctly by all of eight
classical models (two forests, gradient boosting, an MLP, $k$-NN, an SVM, LDA and
logistic regression) and 235 (13\%) by none, and a vote of these cannot repair the
latter. The
forest and the GRU lose their points on the same three gestures but not always through
the same confusion: the GRU sends fist to thumb up 14 points more often, the forest
sends the open hand to the right wave 13 points more often.

\begin{table}[!htb]
  \caption{Share of correctly classified windows per gesture (\%), wearer held out:
  soft vote of ET-500 and RF-500-leaf3 (ML) and GRU (NN), single run of the
  combination test. The last column is the most frequent wrong answer of the ML vote.}
  \label{tab:gestures}
  \centering
  \footnotesize
  \setlength{\tabcolsep}{3.5pt}
  \begin{tabular}{lrrl}
    \hline
    Gesture & ML & NN & Main confusion \\
    \hline
    Rest        & 93.0 & 92.5 & thumb up \\
    Wave left   & 98.4 & 93.3 & -- \\
    Wave right  & 91.6 & 85.2 & open hand \\
    Fist        & 68.3 & 48.7 & open hand \\
    Open hand   & 59.7 & 60.5 & fist, wave right, thumb up \\
    Thumb up    & 59.2 & 54.2 & open hand \\
    \hline
  \end{tabular}
\end{table}

\begin{figure}[!t]
  \centering
  \safefig{width=\columnwidth}{figures/generalisation_gap_paper.png}{generalisation\_gap\_paper.png}
  \caption{The price of an unseen wearer: accuracy per window for a wearer who is in
  the training data (other recordings held out) and for a wearer who is not.}
  \label{fig:gap}
\end{figure}
For a wearer who is in the training data, with other recordings of theirs held out,
both models read 97.4\% (ML vote) and 95.3\% (GRU) of the windows correctly, 20 and 23
points more than for an unseen wearer (Fig.~\ref{fig:gap}); even this figure is
optimistic, because clips cut from one take are close relatives.

Per wearer the picture is uneven (ML vote and GRU, mean of three seeds, each wearer
held out in turn): Betka 82.7\% and 86.0\%, Bruno 79.8\% and 69.2\%, Dano 88.9\% and
74.7\%, Miso 87.7\% and 72.1\%, and Barbora 45.9\% and 36.7\%. The spread between
wearers is about 43 points for the ML vote and 49 for the GRU, and Barbora alone
accounts for most of it. The GRU is better than the vote on one wearer (Betka) and
worse on the other four, but its single seeds differ a great deal (for Barbora 48.8,
29.8 and 31.6\%), and with five wearers we do not read more into this. Why Barbora is hard is not settled: a rotated band does not explain it
(Section~\ref{sec:experiments}), and a loose fit of the band on a thin arm is a
hypothesis that we have not tested.

\subsection{What the models read}
To find out what each model uses, every window was replaced by its eight channel
averages, which removes all shape and timing, and the same model was trained and scored
again (the amplitude-only control). For the ML vote this changes nothing: 77.8\% with the
averages alone against 77.6\% with all 50 level features. The tree ensembles read the
calibrated activation level of each channel, and only that. This is not a defect, since
after per-wearer calibration these levels transfer well between people, but it means
the classical models cannot separate two gestures that differ only in how they
develop in time. The GRU, in contrast, loses 9.2 points under the control (72.4\% to
63.3\%), so it does use the time course of the signal. The loss is not spread evenly:
it is about 26 points for rest, 20 for thumb up and 14 for the open hand, and nothing
for the waves. Two of these three are also the hardest gestures of the ML vote, which
suggests that timing carries some of the information that levels do not. The GRU
result is a single seed, however, and its seed-to-seed spread of about 10 points is
larger than this difference, so we treat it as an indication and not as a measured
effect.

\subsection{How the best models work}
\textbf{Random forest and extra-trees (best classical models).} The best model of the
zoo is a random forest of 500 trees in which every leaf holds at least three windows
(RF-500-leaf3, 80.8\%), followed by the same forest without the leaf limit and by
extra-trees with 500 trees. A forest is a set of decision trees, each trained on a
random part of the windows and a random part of the features, which then vote; extra-trees
additionally choose the split thresholds at random, which makes them faster and
smoother. The input is a vector of 50 amplitude features computed from a 1.0\,s window
of the low-pass filtered, calibrated signal: for each of the eight channels the mean,
median, 90th percentile, standard deviation, the share of the total activation that
falls on the channel and the logarithm of the mean, plus the logarithm of the total
activation and the ratio of the strongest channel to the mean channel. Nothing in this vector describes the shape of
the waveform, which is why the amplitude-only control above loses nothing. The trees
split mainly on which channels are strong and how strong, and the forests fit these
rules without needing a tuned learning rate, which suits a data set this small.

\textbf{GRU (best network).} The network reads a 0.7\,s window, about 40 time steps of
eight channels, after the calibration and a per-channel standardisation fitted on the
training wearers. A single gated recurrent layer with 32 units runs along the window
and produces a state for every time step. Instead of using only the last state, an
attention layer learns a weight for each step and the network keeps both the weighted
mean and the plain mean of the states, so that it can focus on the part of the window in
which the gesture is present without losing the overall activation level. The 64
resulting numbers go through dropout (0.25) and one linear layer to the six classes.
In total it has 4\,455 parameters (25\,KB).

\textbf{TSMixer (second network).} TSMixer is the best network that has no recurrence
and the next one in accuracy (66.3\%, 46\,KB). The eight channels are first embedded
into 24 features per time step. Two mixer blocks then alternate between two small
networks with a residual connection and layer normalisation: one mixes information
along time, with the same weights for every feature, and the other mixes the features of
one time step. The mean and the maximum over time are concatenated and classified by one
linear layer. It reaches almost the accuracy of the GRU while all time steps are
processed in parallel, but it is twice as large and slower per window (1.8 against
0.8\,ms), and in our runs it was not better than the GRU.

\subsection{Speed, soft voting and the single network}
A live prediction is made every 0.15\,s, so everything faster than that keeps up. The
classical pipeline needs 1.7\,ms to resample, calibrate and compute the features and
then about 25\,ms for each of the two forests, so the soft vote of both takes about
52\,ms per window on one CPU thread. The SVM and LDA take under a millisecond, but
they are also 13--14 points less accurate. The network needs about 0.8\,ms for the
forward pass and 1.8\,ms including the round trip to the process that runs it on the
GPU (2.9\,ms on the CPU), about thirty times less than the forest vote, and it is two
orders of magnitude smaller on disk (25\,KB against 5--13\,MB per forest).

Soft voting averages the class probabilities of several models. We searched all
subsets of two to four of the eight best models per family (154 subsets). The best
vote, of ET-500 and RF-500-leaf3, reached 77.6\%, only 0.2 points above ET-500 alone
(77.4\%), and 74.0\% when the search was repeated without each wearer. A vote that
also contained an SVM and LDA was worse (69.5\%), because an unweighted vote gives
the weaker members the same say as the forests. Since two forests also cost twice the
time of one, we regard voting as not worth it here: a single forest is as good and
half as slow. The network is used on its own and is never combined with the classical
models, so that it stays an independent second opinion. Alone it is 5 points behind the
vote in the single-run comparison (72.4\% against 77.6\%) and 13 points behind the best
forest of the zoo when averaged over seeds (67.7\% against 80.8\%), but it is the only
model that uses the time course of the signal, it is tiny, and it is fast enough to run
on a small board. The gap between the two numbers for the GRU is mainly seed
variance.

\subsection{Limitations}
The results rest on five wearers, and the figures should be read with that in mind.
Each wearer has only 17 to 102 windows of a gesture, from a handful of recordings, so
a single cell of a per-wearer table can move a lot with one more recording, and
the ML score depends strongly on one person (Barbora). The smallest cell, Betka's
thumb up (17 windows), is small for a recorded reason: the first seven of her 14
thumb-up clips (3\,s each) were discarded and seven shorter clips (1.0--2.1\,s) were
recorded afterwards, and only these are in the dataset. The soft vote of the two forests
reached 77.6\% on the folds on which it was selected, but 74.0\% when the whole
selection was repeated without each wearer, so the first number is optimistic. The
network figures come from one or two training seeds and differ by up to about 10
points between seeds, which is more than most of the differences between networks.

The electrode contact is another limitation. The band does not fit all arms equally
tightly, and one of the wearers, who has a thin arm and is the hardest case, was observed by
the operator not to wear it tightly; the fit was not measured. We did not record a
second sitting with a tighter fit, so whether a loose contact is the cause remains
untested. Calibration also has a cost: a model trained on
calibrated features needs every new wearer to calibrate first, and the calibration used
for training was derived from the recorded signal rather than from the live
rest-and-squeeze procedure. Only one wearer (Dano) has a live calibration, a guided
sweep over the gestures, and its per-channel reference levels differ from the derived
ones by a factor of 0.8 to 2.2; whether a model trained on the derived scale predicts
well on the live scale was not measured, and the other wearers have no live
calibration.

Finally, all accuracies are measured offline on recorded holds, with the wearer
held out, and not live. Live use differs in three ways: gestures are made and relaxed
rather than held, the prediction is repeated every 0.15\,s on overlapping windows, and
the wearer sees the result. The live accuracy has not been measured, and the numbers
in this paper should not be read as one.
\textcolor{red}{[TODO: measure the live accuracy in a real session with the armband worn,
not on the recorded data set, and report it here.]}

\section{Conclusion}
\label{sec:conclusion}
We compared classical machine learning with small neural networks for recognising six
gestures from the Myo armband on five wearers, always on a wearer the model had not
seen. The best classical model, a random forest, reached 80.8\% per window with
calibration and adaptation to the wearer; the soft vote of two forests reached 77.6\%
(74.0\% when the selection was repeated without each wearer) and the best network, a GRU
with 4\,455 parameters, 72.4\% in a single run and 67.7\% averaged over seeds. The
classical models are therefore ahead by about 5 to 13 points, but the network is two
orders of magnitude smaller and about thirty times faster per window. Voting did not
help, since one forest is as good as two.

Calibrating every wearer on their own signal was the most useful addition, and adapting
the input statistics to the wearer helped the classical models further. Augmentation did
not help. The forests read only the calibrated level of each channel, whereas the GRU
also uses the time course of the signal, which is the reason to expect it to improve with
more data. Most errors come from fist, open hand and thumb up, and from one wearer
(Barbora), whose recordings the models read at only about 60\% (ML) and 39\% (networks)
when they have not seen her. With five wearers these results are an indication and
not a final ranking.

The next steps are more wearers, a second sitting with a tighter fit of the band for
the hardest wearer, and a measurement of the live accuracy.
\textcolor{red}{[TODO: live measurement on a real armband, not from the recorded data
set, still has to be done and added here.]}
The code, the recording tool and the training pipeline are available at
\textcolor{red}{\texttt{https://github.com/Inszert/emg}}
\textcolor{red}{[TODO: replace with the official repository.]}

\printbibliography

\end{document}